\section{Introduction}

The ideal intersection graph associates a graph to the intersection
structure of the nonzero proper ideals of a ring~\cite{intersection}.
For $\mathbb Z_n$, its isomorphism type depends on the exponents in the
prime factorization of $n$. We ask which of these graphs have integer
Laplacian eigenvalues.

We prove a complete classification for squarefree composite integers
(Theorem~\ref{thm:squarefree}), an explicit quadratic obstruction for
$(1,1,k)$ (Theorem~\ref{thm:pqrk}), and a linear cutoff for repeated
exponents (Theorem~\ref{thm:second}). The cutoff combines an arithmetic
restriction on the antisymmetric block with the symmetric cubic obstruction
for $b=(a-1)(2a-1)$ proved in Theorem~\ref{thm:boundary}.
An exact divisor criterion isolates every possible square discriminant for
fixed $a$, leading to nonintegrality for all $b$ when $a=2,3,4,5,6$
(Theorem~\ref{thm:small} and Corollary~\ref{cor:five-six}).

The graph and lifting arguments are given explicitly, so the proofs do not
require a complete spectral decomposition. Section~\ref{sec:verification}
describes the exact computations supporting the displayed formulas and
finite certificates. The classification for arbitrary exponent vectors
remains open.
